\clearpage
\begingroup
\tcbset{myboxstyle/.append style={breakable=false}}
\section{Математическое ожидание и дисперсия}

\subsection{Математическое ожидание}

Пусть $(\Omega,\mathcal F,\Pr)$ --- вероятностное пространство,
а $\xi\colon\Omega\to\R$ --- случайная величина.
Её распределение задаётся вероятностной мерой
\[
    \Pr_\xi(B)=\Pr\{\xi\in B\},\qquad B\in\mathcal B(\R).
\]

\subsubsection{Конечная модель и среднее арифметическое}

Пусть $\Omega=\{\omega_1,\ldots,\omega_n\}$. Проведём $N$ опытов
с результатами $\omega^1,\ldots,\omega^N$ и обозначим через $N_i$
число появлений исхода $\omega_i$. Тогда среднее наблюдаемых значений равно
\[
    \frac1N\sum_{j=1}^{N}\xi(\omega^j)
    =\sum_{i=1}^{n}\xi(\omega_i)\frac{N_i}{N}.
\]
В вероятностной модели частотам $N_i/N$ соответствуют вероятности
$\Pr\{\omega_i\}$.

\begin{definition}[Математическое ожидание в конечной модели]
    \textbf{Математическим ожиданием} случайной величины $\xi$ называется
    \[
        \EE\xi=\sum_{i=1}^{n}\xi(\omega_i)\Pr\{\omega_i\}.
    \]
    Используется также обозначение $M\xi$.
\end{definition}

\begin{definition}[Общее определение]
    Математическое ожидание случайной величины $\xi$ --- интеграл Лебега
    \[
        \EE\xi=\int_\Omega\xi(\omega)\,\diff\Pr(\omega).
    \]
    Конечное математическое ожидание существует при условии
    \[
        \int_\Omega|\xi(\omega)|\,\diff\Pr(\omega)<\infty.
    \]
\end{definition}

\begin{theorem}[Вычисление по распределению]
    Если в конечной модели $\xi$ принимает различные значения
    $x_1,\ldots,x_k$ с вероятностями $p_i=\Pr\{\xi=x_i\}$, то
    \[
        \EE\xi=\sum_{i=1}^{k}x_i p_i.
    \]
\end{theorem}

\begin{Proof}
    Сгруппируем исходы, на которых $\xi$ принимает одинаковое значение:
    \[
        \sum_{\omega\in\Omega}\xi(\omega)\Pr\{\omega\}
        =\sum_{i=1}^{k}x_i
          \sum_{\omega:\,\xi(\omega)=x_i}\Pr\{\omega\}
        =\sum_{i=1}^{k}x_i p_i.
    \]
\end{Proof}

\newpage
\subsubsection{Дискретное и абсолютно непрерывное распределения}

\begin{definition}[Дискретное распределение]
    Для распределения $\{(x_i,p_i)\}_i$ математическое ожидание равно
    \[
        \EE\xi=\sum_i x_i p_i,
        \qquad\text{если}\quad \sum_i |x_i|p_i<\infty.
    \]
    Таким образом, требуется абсолютная сходимость ряда.
\end{definition}

\begin{example}[Распределение Пуассона]
    Если $\xi\sim\operatorname{Poiss}(\lambda)$, $\lambda>0$, то
    \[
        \Pr\{\xi=k\}=e^{-\lambda}\frac{\lambda^k}{k!},
        \qquad k=0,1,2,\ldots
    \]
    Все слагаемые неотрицательны, и
    \[
        \begin{aligned}
            \EE\xi
            &=\sum_{k=0}^{\infty}k e^{-\lambda}\frac{\lambda^k}{k!}
             =\lambda e^{-\lambda}
                \sum_{k=1}^{\infty}\frac{\lambda^{k-1}}{(k-1)!}\\
            &=\lambda e^{-\lambda}
                \sum_{m=0}^{\infty}\frac{\lambda^m}{m!}
             =\lambda e^{-\lambda}e^\lambda=\lambda.
        \end{aligned}
    \]
\end{example}

\begin{definition}[Абсолютно непрерывное распределение]
    Если $\xi$ имеет плотность $p_\xi$, то
    \[
        \EE\xi=\int_{-\infty}^{+\infty}x p_\xi(x)\,\diff x,
        \qquad\text{если}\quad
        \int_{-\infty}^{+\infty}|x|p_\xi(x)\,\diff x<\infty.
    \]
    Если интеграл записан как несобственный интеграл Римана,
    требуется его абсолютная сходимость.
\end{definition}

Для малого промежутка $\Delta_i$ около точки $x_i$, в которой плотность
непрерывна,
\[
    \Pr\{\xi\in\Delta_i\}\approx p_\xi(x_i)|\Delta_i|.
\]
Поэтому суммам $\sum_i x_i\Pr\{\xi\in\Delta_i\}$ соответствуют
интегральные суммы $\sum_i x_i p_\xi(x_i)|\Delta_i|$.

\begin{consequence}
    Математическое ожидание определяется распределением случайной величины:
    \[
        \EE\xi=\int_\R x\,\diff\Pr_\xi(x).
    \]
\end{consequence}

\newpage
\subsection{Свойства математического ожидания}

\begin{theorem}[Ожидание константы и линейность]
    Для $c\in\R$ имеем $\EE c=c$.
    Если существуют конечные $\EE\xi_1$ и $\EE\xi_2$, то для любых
    $a_1,a_2\in\R$ существует математическое ожидание их линейной комбинации и
    \[
        \EE(a_1\xi_1+a_2\xi_2)
        =a_1\EE\xi_1+a_2\EE\xi_2.
    \]
\end{theorem}

\begin{Proof}[Линейность в дискретном случае]
    Пусть $\xi_1$ принимает значения $x_i$, а $\xi_2$ --- значения $y_j$.
    Обозначим
    \[
        p_{ij}=\Pr\{\xi_1=x_i,\ \xi_2=y_j\},\qquad
        p_i=\Pr\{\xi_1=x_i\},\quad q_j=\Pr\{\xi_2=y_j\}.
    \]
    По счётной аддитивности вероятности
    \[
        \sum_j p_{ij}=p_i,\qquad \sum_i p_{ij}=q_j.
    \]
    Сначала проверим абсолютную сходимость:
    \[
        \begin{aligned}
            \sum_{i,j}|a_1x_i+a_2y_j|p_{ij}
            &\leq |a_1|\sum_i|x_i|\sum_jp_{ij}
                  +|a_2|\sum_j|y_j|\sum_ip_{ij}\\
            &=|a_1|\sum_i|x_i|p_i+|a_2|\sum_j|y_j|q_j<\infty.
        \end{aligned}
    \]
    Следовательно, суммы можно перегруппировать:
    \[
        \begin{aligned}
            \EE(a_1\xi_1+a_2\xi_2)
            &=\sum_{i,j}(a_1x_i+a_2y_j)p_{ij}\\
            &=a_1\sum_i x_i p_i+a_2\sum_j y_j q_j\\
            &=a_1\EE\xi_1+a_2\EE\xi_2.
        \end{aligned}
    \]
\end{Proof}

\begin{theorem}[Монотонность]
    Если $\xi\geq\eta$ почти наверное и оба математических ожидания конечны,
    то $\EE\xi\geq\EE\eta$.
\end{theorem}

\begin{Proof}
    Для $\zeta=\xi-\eta\geq0$ имеем $\EE\zeta\geq0$.
    По линейности $\EE\zeta=\EE\xi-\EE\eta$, откуда
    $\EE\xi\geq\EE\eta$.
\end{Proof}

\newpage
\begin{theorem}[Оценка модуля ожидания]
    Если $\EE\xi$ конечно, то
    \[
        |\EE\xi|\leq\EE|\xi|.
    \]
\end{theorem}

\begin{Proof}
    Ожидание $\EE|\xi|$ конечно по условию абсолютной интегрируемости.
    Из $-|\xi|\leq\xi\leq|\xi|$, монотонности и линейности следует
    \[
        -\EE|\xi|\leq\EE\xi\leq\EE|\xi|.
    \]
\end{Proof}

\subsection{Дисперсия}

\begin{example}[Одинаковое ожидание и различный разброс]
    Пусть
    \[
        \Pr\{\xi_1=-1\}=\Pr\{\xi_1=1\}=\frac12,\qquad
        \Pr\{\xi_2=-100\}=\Pr\{\xi_2=100\}=\frac12.
    \]
    Тогда $\EE\xi_1=\EE\xi_2=0$, но значения второй величины
    дальше от математического ожидания.
\end{example}

\begin{definition}
    \textbf{Дисперсией} случайной величины $\xi$ с конечным математическим
    ожиданием называется
    \[
        \DD\xi=\operatorname{Var}\xi
        =\EE(\xi-\EE\xi)^2.
    \]
    Конечная дисперсия существует, если $\EE\xi^2<\infty$.
\end{definition}

\begin{theorem}[Формула для вычисления дисперсии]
    При $\EE\xi^2<\infty$ справедливо равенство
    \[
        \DD\xi=\EE\xi^2-(\EE\xi)^2.
    \]
\end{theorem}

\begin{Proof}
    По линейности математического ожидания
    \[
        \begin{aligned}
            \DD\xi
            &=\EE\bigl(\xi^2-2\xi\EE\xi+(\EE\xi)^2\bigr)\\
            &=\EE\xi^2-2(\EE\xi)^2+(\EE\xi)^2\\
            &=\EE\xi^2-(\EE\xi)^2.
        \end{aligned}
    \]
\end{Proof}
\endgroup
